\chapter{Theoretical Background}
\label{ch:theory}

This chapter provides the mathematical background for the methods used in this work. The
techniques are presented in the order in which the methodology applies them: first the
matrix factorization and clustering tools used to separate tumor spectra from background
(Sections~\ref{sec:nmf} and~\ref{sec:gmm}), then the neural network theory underlying the
generative model (Sections~\ref{sec:deep_learning} and~\ref{sec:cvae}), then the downstream
classifiers evaluated in Chapter~\ref{sec:results} (Section~\ref{sec:classifiers}), and
finally the embedding and statistical tools used to analyse the results
(Sections~\ref{sec:embeddings} and~\ref{sec:airpls_bootstrap}).


\section{Non-Negative Matrix Factorization}
\label{sec:nmf}

We will present here the problem of Non-Negative Matrix Factorization (NMF) with an emphasis on hyperspectral decomposition applications. This is a generic optimization problem with many applications and various theoretical formulations can be found, but we will formulate the problem in our context, as presented by Gillies \cite{Gillis2017,Gillis2020}.

The NMF can be formulated as follows: given a non-negative matrix $M \in \mathbb{R_+^{m \times n}}$, a factorization rank $r$, and a distance measure $D(\cdot,\cdot)$ find two non-negative matrices $W \in \mathbb{R}_+^{m \times r}$ and $H \in \mathbb{R}_+^{r \times n}$ such that $D(M,WH)$ is minimized, i.e. solve:
$$ \min_{W \in \mathbb{R}_+^{m \times r}, H \in \mathbb{R}_+^{r \times n}} D(M,WH)$$

In practice, we will consider the distance $D$ to be the Frobenius norm of the error matrix $M-WH$, hence the problem can be writen:

$$ \min_{W \in \mathbb{R}_+^{m \times r}, H \in \mathbb{R}_+^{r \times n}} \norm{M-WH}_F^2$$

This is the most natural and common choice as it corresponds to the assumption of i.i.d. Gaussian noise and it leads to a smooth optimization problem which is easy to solve \cite{Gillis2020}

There are many variants of NMF, for example the factorization can be solved exactly, in which case the problem is called Exact NMF, but in most applications we are only interested in an approximate solution. Also, for most solutions, additional constraints can be imposed.

The non-negativity constraints lead to natural physical interpretations of the terms of the factorization. The intuition is that spectral signatures are additive combinations of spectral signatures of the combining elements. For example, in our setting of biological samples, the spectral signature of a probe could be a linear combination of spectral signatures of lipids, proteins, phenilalanine, DNA C-H molecules, etc. This NMF factorization is called endmember unmixing, or more precisely, blind endmember unmixing, where individual components are called endmembers. Due to the additive nature of the signal acquisition process, where Raman bands appear where there is a certain molecule that has an absorbtion energy band of a specific energy level, this formulation of the NMF is well suited for the physical process involved.

\subsection*{Sparsity of factorization}

Non-negativity constraints are linear inequality constraints that naturally enforce sparsity, or solutions that contain lots of zeros. This derives from the first order optimality conditions (Karush-Kuhn-Tucker) of an optimization problem with such constraints \cite{Kuhn1951,Gillis2020}:

$$ \min_{x \in \mathbf{R^n}} f(x) \text{such that} x \ge 0 $$
giving the conditions:
$$ x \ge 0, \nabla f(x) \ge 0, \text{and} x_i(\nabla f(x))_i=0 \text{for all i.}$$

This enforces $x_i = 0$ if $(\nabla f(x))_i > 0$, which leads to sparse solutions.

\subsection*{Solution and Convergence}
Many algorithms have been developed for solving the NMF problem \cite{ChihJenLin2007,Xu2016,Lee2000}. We present here a simple solution by Daniel D. Lee and H. Sebastian Seung \cite{Lee2000}.

\begin{theorem}
    The distance $\norm{M-WH}_F^2$ is nonincreasing under the update rules
    $$ H_{a \mu} \leftarrow H_{a \mu} \frac{(W^T M)_{a \mu}}{(W^T W H)_{a \mu}} \qquad
    W_{i a} \leftarrow W_{i a} \frac{(M H^T)_{i a}}{(W H H^T)_{i a}}$$
    This distance is invariant under these updates iff $W$ and $H$ are at a stationary point.
\end{theorem}

The proof of convergence can be found in full in D.D. Lee's article. We can see that the updates alternate, hence it resembles the Expectatio-Maximization algorithm. \cite{Lee2000}.

In our analysis we used the scikit-learn implementation of the NMF \cite{sklearn2011,Cichocki2009}.

\section{Gaussian Mixture Model}
\label{sec:gmm}

The Gaussian Mixture Model (GMM) is a widely used tool that assigns soft labels to a set of data points. We can consider it a soft clustering technique, where we define a hard clustering as a partition of the set of points and a soft clustering a labeling of the points with a probability distribution over the number of clusters.
We present here a short description of the method as described by C.H. Bishop in "Pattern Recognition and Machine Learning", Chapter 9 \cite{bishop2006pattern}.

The Gaussian Mixture can be defined as a weighted sum of Gaussian distributions:

$$ p(\mathbf{x}) = \sum_{k=1}^{K}\pi_k \mathcal{N}(\mathbf{x} | \mu_k, \Sigma_k) $$

where $\mathbf{x}$ is the data set $\mathbf{x} = \{x_1, x_2, \dots x_N\}$. In order to assign points to clusters, we will assume there exists this assignment as a latent variable. We will denote this latent variariable with $\mathbf{z}$, as unit vectors of $\mathbb{R}^K$. We can now define the joint distribution $\mathbb{P}(\mathbf{x},\mathbf{z})$, the marginal distribution $\mathbb{P}(\mathbf{z})$, the conditional distribution $\mathbb{P}(\mathbf{x}|\mathbf{z})$. The marginal distribution over $\mathbf{z}$ is:
$$\mathbb{P}(z_k = 1) = \pi_k$$
where the parameters $\{\pi_k\}$ form a probability distribution, so they must satisfy the constraints:
$$0 \le \pi_k \le 1 \quad \sum_{k=1}^{K}\pi_k = 1$$
Because $z_k \in \{0,1\}$, we can write:
$$\mathbb{P}(\mathbf{z}) = \prod_{k=1}^K \pi_k^{z_k}$$
We make the assumption that the conditional distribution of x given z is a Gaussian:
$$\mathbb{P}(\mathbf{x}|z_k = 1) = \mathcal{N}(\mathbf{x} | \mu_k, \Sigma_k)$$

By summing over all states of z we obtain the marginal distribution of x:

$$ \mathbb{P}(x) = \sum_{\mathbf{z}} \mathbf{P}(\mathbf{z})\mathbf{P}(\mathbf{x}|\mathbf{z}) = \sum_{k=1}^{K} \pi_k \mathcal{N(\mathbf{x}|\mu_k, \Sigma_k)}$$

Of particular interest is also the conditional probability of $\mathbf{z}$ given $\mathbf{x}$. We will use the following notation:
$$\gamma(z_k) := \mathbb{P}(z_k=1|x) = \frac{\mathbb{P}(z_k = 1) \mathbb{P}(\mathbf{x}|z_k = 1) }{\sum_{j=1}^{K} \mathbb{P}(z_j = 1) \mathbb{P}(\mathbf{x}|z_j = 1} = \frac{\pi_k \mathcal{N}(\mathbf{x}|\mu_k, \Sigma_k)}{\sum_{j=1}^{K} \pi_j \mathcal{N}(\mathbf{x}|\mu_j, \Sigma_j)}$$

We can now view $\pi_k$ as the prior probability and $\gamma(z_k)$ as the posterior after observing $x$.

With $\mathbf{x}$ as our dataset, the log of the likelihood can be written as:
$$ \ln \mathbb{P}(\mathbf{x}|\pi,\mu,\Sigma) = \sum_{n=1}^{N} \ln \left\{\sum_{k=1}^{K} \pi_k \mathcal{N}(\mathbf{x}_n|\mu_k,\Sigma_k)\right\} $$

To find the maximum of this function, we set the derivatives with respect to $\mu_k$ to 0 and we obtain:
$$ \mu_k = \frac{1}{N_k} \sum_{n=1}^{N}\gamma(z_{nk})\mathbf{x}_n $$
where
$$ N_k = \sum_{n=1}^{N} \gamma(z_{nk}) $$

Next we set the derivatives with respect to $\Sigma_k$ to zero and we obtain:

$$ \Sigma_k = \frac{1}{N_k} \sum_{n=1}^{N}\gamma(z_{nk})(\mathbf{x}_n - \mu_k)(\mathbf{x}_n - \mu_k)^T $$

In order to find the $\pi_k$ coeficients, we take into account the probability constraints over $\pi_k$ and we maximize:
$$ \ln \mathbb{P}(\mathbf{x}|\pi, \mu, \Sigma) + \lambda \left(\sum_{k=1}^{K} \pi_k - 1 \right)$$
and we obtain
$$ \pi_k = \frac{N_k}{N} $$

\subsection*{The Expectation Maximization algorithm}

We cannot find a direct solution to the GMM problem, but we can devise an iterative algorithm by alternating between the cluster assignments and the parameters of the gaussians. The algorithm is called Expectation Maximization (EM):



\begin{algorithm}
\caption{The EM Algorithm}
\begin{algorithmic}[1] 
    \State \textbf{Initialize:} Randomly initialize $\mu_k$, $\Sigma_k$, $\pi_k$ 

    
    \While{not converged} 
        
        \State \textbf{E step: evaluate the clustering}
        \vspace{1.5ex} % Spacing before the E-step formula
        \Statex \hspace{\algorithmicindent} $\displaystyle \gamma(z_{nk}) = \frac{\pi_k \mathcal{N}(\mathbf{x}_n|\mu_k,\Sigma_k)}{\sum_{j=1}^{K} \pi_j \mathcal{N}(\mathbf{x}_n|\mu_j,\Sigma_j)}$
        \vspace{2.5ex} % Spacing after the E-step formula
        
        \State \textbf{M step: estimate parameters using the current clustering}
        \vspace{1.5ex} % Spacing before the M-step aligned formulas
        \Statex \hspace{\algorithmicindent} $\begin{aligned}
            \mu_k^{\text{new}} &= \frac{1}{N_k} \sum_{n=1}^{N} \gamma(z_{nk}) \mathbf{x}_n \\[1.5ex]
            \Sigma_k^{\text{new}} &= \frac{1}{N_k} \sum_{n=1}^{N} \gamma (z_{nk})(\mathbf{x}_n - \mu_k^{\text{new}})(\mathbf{x}_n - \mu_k^{\text{new}})^T \\[1.5ex]
            \pi_k^\text{new} &= \frac{N_k}{N} \quad \text{where} \quad N_k = \sum_{n=1}^{N} \gamma(z_{nk})
        \end{aligned}$
        \vspace{2.5ex} % Spacing after the M-step formulas
        
        \State \textbf{Evaluate the log likelihood for the entire dataset} 
        \vspace{1.5ex} % Spacing before the log-likelihood formula
        \Statex \hspace{\algorithmicindent} $\displaystyle \ln \mathbb{P}(\mathbf{x}|\mu, \Sigma, \pi) = \sum_{n=1}^{N} \ln \left( \sum_{k=1}^{K} \pi_k \mathcal{N}(\mathbf{x}_n| \mu_k, \Sigma_k) \right)$
        \vspace{2ex} % Spacing before the loop ends

    \EndWhile
\end{algorithmic}
\end{algorithm}

Note that this is an iterative algorithm and it converges to a local solution. While this works in practice most of the time, there exists the posibility that the solution is not optimal.

\section{Deep Learning}
\label{sec:deep_learning}

Research in Artificial Neural Networks, or Deep Learning has expanded significantly in the last decades, motivated by advances in computing power, development of software frameworks for easy implementation and testing, as well as theoretical results that offer some guarantees regarding the behaviour of Artificial Neural Networks. 

\subsection*{Early results: ANNs as Universal Function Approximators}
The first theoretical results were formulated in 1989, with Kurt Hornik et al. \cite{Hornik1989} prooving that a Neural Network with one hidden layer, with a squashing activation function can aproximate arbitrarily well any continuous function on a compact set, provided we let the number of neurons in the hidden layer grow arbitrarily large, and G. Cyabenko proving the Universal Approximator theorem for sigmoid activation functions \cite{Cybenko1989}.

\subsection*{Extending the Universal Function Approximation Theorem}
Other activation functions were tested, and as the requrements on the function were reduced to continuous, non linear, and differentiable almost everywhere, the function $X^+$ (0 if $x < 0$, $x$ otherwise), known in the Deep Learning comunity as the Rectified Linear Unit (ReLU) activation function became widespread.

The theorem of Universal Approximation was proved for the ReLU activation function as well under some particular geometries \cite{Hanin2019, Lu2017}. 

Other variants of the theorem relevant for our design are the ones regarding the convolution operation \cite{Zhou2020, Yarotsky2018}. This acts as a shared local operator replacing the full matrix multiplication .

We observe that there are many variants of the Universal Approximation Theorem, and many different proofs for various ANNs different geometries, activation function or other components like local operators.


\section{Conditional Variational Autoencoder}
\label{sec:cvae}

The generative model used in this work is a Conditional Variational Autoencoder (CVAE). This
section develops it in three steps: the variational autoencoder as a latent variable model
\cite{kingma2014vae}, the Evidence Lower Bound that makes it trainable, and the conditioning
mechanism that turns it into a class-conditional generator \cite{sohn2015cvae}. The
architecture and the training objective actually used are given in
Section~\ref{sec:cvae_architecture}.

\subsection*{The latent variable model}

We assume the observed spectra $x \in \mathbb{R}^{N_\lambda}$ are generated by a two-stage
process: a latent variable $z \in \mathbb{R}^{d}$ is drawn from a prior $p(z)$, and the
spectrum is then drawn from a conditional distribution $p_\theta(x|z)$ parameterised by a
neural network with parameters $\theta$. The marginal likelihood of a spectrum is therefore
\begin{equation}
p_\theta(x) = \int p_\theta(x|z)\, p(z) \, dz .
\label{eq:marginal_likelihood}
\end{equation}

Maximising \eqref{eq:marginal_likelihood} directly is intractable: the integral runs over the
whole $d$-dimensional latent space, and $p_\theta(x|z)$ is a neural network, so no closed form
exists. The same intractability affects the true posterior
\begin{equation}
p_\theta(z|x) = \frac{p_\theta(x|z)\,p(z)}{p_\theta(x)},
\label{eq:true_posterior}
\end{equation}
whose denominator is exactly the intractable quantity \eqref{eq:marginal_likelihood}.

The variational approach sidesteps this by introducing a second neural network, the
\emph{encoder} $q_\phi(z|x)$ with parameters $\phi$, which approximates the true posterior
$p_\theta(z|x)$. Rezende et al.\ \cite{rezende2014stochastic} show that parameterising the
mean and variance of a Gaussian by neural networks in this way is sufficient to approximate
complex posterior densities, which justifies the Gaussian choice made below.

\TODO{One paragraph of intuition here: the encoder compresses a spectrum into a
low-dimensional code and the decoder reconstructs it, but unlike a plain autoencoder both
maps are distributions rather than points. Say why a distribution is what makes sampling
possible later.}

\subsection*{The Kullback--Leibler divergence}

For two distributions $q$ and $p$ over the same space, the Kullback--Leibler divergence is
\begin{equation}
D_{KL}(q \| p) = \mathbb{E}_{q}\left[ \log \frac{q(z)}{p(z)} \right]
              = \int q(z) \log \frac{q(z)}{p(z)} \, dz .
\label{eq:kl_definition}
\end{equation}
It is non-negative, and zero if and only if $q = p$ almost everywhere. It is not symmetric, so
it is not a metric; $D_{KL}(q\|p)$ measures the cost of describing samples from $q$ using a
code built for $p$.

\subsection*{The Evidence Lower Bound}

Starting from the log marginal likelihood and introducing $q_\phi(z|x)$:
\begin{align}
\log p_\theta(x)
  &= \log \int p_\theta(x|z)\,p(z)\,dz
   = \log \int q_\phi(z|x) \frac{p_\theta(x|z)\,p(z)}{q_\phi(z|x)} \, dz \nonumber \\
  &= \log \mathbb{E}_{q_\phi(z|x)}\!\left[ \frac{p_\theta(x|z)\,p(z)}{q_\phi(z|x)} \right]
   \;\ge\; \mathbb{E}_{q_\phi(z|x)}\!\left[ \log \frac{p_\theta(x|z)\,p(z)}{q_\phi(z|x)} \right],
\label{eq:jensen}
\end{align}
where the inequality is Jensen's inequality applied to the concave function $\log$. Splitting
the logarithm gives the Evidence Lower Bound (ELBO):
\begin{equation}
\mathcal{L}_{\text{ELBO}}(\theta,\phi;x)
 = \underbrace{\mathbb{E}_{q_\phi(z|x)}\big[ \log p_\theta(x|z) \big]}_{\text{reconstruction}}
 - \underbrace{D_{KL}\big(q_\phi(z|x) \,\|\, p(z)\big)}_{\text{regularisation}} .
\label{eq:elbo}
\end{equation}

The gap between the bound and the true log-likelihood is itself a KL divergence,
\begin{equation}
\log p_\theta(x) - \mathcal{L}_{\text{ELBO}}(\theta,\phi;x)
 = D_{KL}\big(q_\phi(z|x) \,\|\, p_\theta(z|x)\big) \;\ge\; 0 ,
\label{eq:elbo_gap}
\end{equation}
so maximising the ELBO simultaneously raises a lower bound on the evidence and drives the
approximate posterior towards the true one.

\TODO{Two or three sentences reading equation \eqref{eq:elbo} back in words: the first term
rewards faithful reconstruction, the second penalises latent codes that stray from the prior,
and training balances the two. Connect this to the composite loss in
Section~\ref{sec:cvae_architecture}, where the reconstruction term is realised as MSE and
cosine loss and the KL term carries the weight $\beta$.}

\subsection*{The Gaussian case and the reparameterisation trick}

We take the prior to be a standard normal, $p(z) = \mathcal{N}(0, I)$, and the approximate
posterior to be a diagonal Gaussian whose parameters are produced by the encoder,
\begin{equation}
q_\phi(z|x) = \mathcal{N}\big(z \,;\, \mu_\phi(x), \, \operatorname{diag}(\sigma^2_\phi(x))\big).
\label{eq:gaussian_posterior}
\end{equation}
Under these choices the KL term of \eqref{eq:elbo} has the closed form used in
Section~\ref{sec:cvae_architecture}:
\begin{equation}
D_{KL}\big(q_\phi(z|x) \| p(z)\big)
 = \frac{1}{2} \sum_{i=1}^{d} \left( \mu_i^2 + \sigma_i^2 - \log(\sigma_i^2) - 1 \right).
\label{eq:kl_gaussian}
\end{equation}
In practice the encoder emits $\log(\sigma_i^2)$ rather than $\sigma_i$, which keeps the
variance positive without a constraint and is numerically better behaved.

The reconstruction term of \eqref{eq:elbo} is an expectation over $q_\phi(z|x)$, estimated in
training by sampling $z$. Sampling is not differentiable with respect to $\phi$, so the
gradient cannot be propagated back into the encoder. The \emph{reparameterisation trick}
\cite{kingma2014vae} resolves this by moving the randomness into a parameter-free variable:
\begin{equation}
z = \mu_\phi(x) + \sigma_\phi(x) \odot \epsilon,
\qquad \epsilon \sim \mathcal{N}(0, I),
\label{eq:reparameterisation}
\end{equation}
where $\odot$ is the elementwise product. The sample $z$ is now a deterministic, differentiable
function of $\mu_\phi(x)$ and $\sigma_\phi(x)$, and the gradient flows through it while the
stochasticity is confined to $\epsilon$.

\subsection*{Why the KL term turns the decoder into a generator}

\TODO{This subsection is the conceptual keystone of the thesis and deserves a careful
paragraph in your own words. The argument to make: to generate a new spectrum we sample
$z \sim \mathcal{N}(0,I)$ and pass it through the decoder, but this is only meaningful if the
decoder has actually been trained on inputs distributed that way. A plain autoencoder places
its codes wherever reconstruction is easiest, leaving regions of latent space the decoder has
never seen, so a prior sample from such a region decodes to noise. The KL term in
equation \eqref{eq:elbo} penalises exactly this: it pulls the aggregate distribution of
encoded spectra towards the prior, so that the region the decoder was trained on and the
region we sample from coincide. Then state the trade-off: raising $\beta$ tightens that match
at the cost of reconstruction fidelity \cite{higgins2017betavae}, and note that
Chapter~\ref{sec:discussion} finds the resulting synthetic spectra to have visibly narrower
variance than the real ones -- which is this trade-off made visible.}

\subsection*{Conditioning: from VAE to CVAE}

To control which class a generated spectrum belongs to, both encoder and decoder receive the
class label $y$ as an additional input \cite{sohn2015cvae}. The ELBO becomes
\begin{equation}
\mathcal{L}_{\text{ELBO}}(\theta,\phi;x,y)
 = \mathbb{E}_{q_\phi(z|x,y)}\big[ \log p_\theta(x|z,y) \big]
 - D_{KL}\big(q_\phi(z|x,y) \,\|\, p(z|y)\big),
\label{eq:celbo}
\end{equation}
and the prior is usually taken to be independent of the label, $p(z|y) = p(z) = \mathcal{N}(0,I)$,
so that a single latent space is shared across classes and the label alone selects the class.

Generation is then immediate: draw $z \sim \mathcal{N}(0, I)$, choose a one-hot label $y_c$,
and evaluate $\hat{x} = p_\theta(x | z, y_c)$. This is the sampling procedure used to build
the synthetic training sets of Chapter~\ref{sec:experimental_design}, and, run once per
candidate class, it is also the mechanism behind the Classification by Reconstruction strategy
of Section~\ref{sec:cbr}.

\TODO{Optional closing paragraph: note that conditioning the decoder does not by itself
guarantee class-discriminative latent codes, which is the motivation for the auxiliary
classification head described in Section~\ref{sec:cvae_architecture}. Related applications of
conditional VAEs to spectral and time-series data: \texttt{Gebran2025} (stellar spectra),
\texttt{Kavran2022} (beta-VAE augmentation for time series).}
